\section{Metodología}

Para el desarrollo del trabajo se ha seguido la metodología de desarrollo de software en cascada, que se compone de las siguientes fases: análisis, diseño, implementación y validación.

\section{Fases de desarrollo}

\subsection{Análisis}

En esta primera fase se han identificado los requisitos del sistema y se han estudiado los protocolos de comunicación empleados en los vehículos del grupo VAG (Volkswagen Audi Group) sobre la plataforma PQ25, para poder empezar con el diseño de la infraestructura de monitorización y comunicación con las distintas ECUs del vehículo.

Durante este análisis se ha descubierto que los vehículos del grupo VAG sobre la plataforma PQ25 utilizan una topología de \textbf{Gateway Central} esto lo que consigue es diferenciar los buses CAN de forma que no haya saturación de los datos en un solo CAN Bus. Esta arquitectura implica además que el bus CAN-Powertrain \textbf{no es directamente accesible desde el conector OBD-II}: el Gateway actúa como intermediario, filtrando y enrutando los mensajes entre buses, lo que impide el acceso directo al bus de motor mediante un adaptador CAN estándar conectado al puerto de diagnóstico. Podemos diferenciar entre:
\begin{itemize}
    \item \textbf{CAN-Powertrain}: este bus conecta la ECU (Motor), sensor ABS/ESP y transmisión. Este es el CAN Bus más importante con lo que tiene una velocidad de 500 kbps.
    \item \textbf{CAN-Diagnostic}: este es el bus de diagnositco que se revela en el puerto OBD-II situado en la parte inferior del salpicadero bajo el volante. Este es un bus que opera a 500 kbps.
    \item \textbf{CAN-Comfort}: este bus al no ser de alta prioridad tiene una velocidad de 100 kbps y controla elevalunas, cierre centralizado, iluminación interna, etc.
    \item \textbf{CAN-Infotainment}: este bus es el que controla los sistema de entretenimiento e información como la radio, navegación y bluetooth. Este bus también opera a 100 kbps.
\end{itemize}

En cuanto a la capa de Enlace y Red se encontró que se trabaja siguiendo las normas ISO(Organización Internacional de Normalización):
\begin{itemize}
    \item \textbf{ISO 11898-1}: este norma define el estándar físico y de enlace de datos del CAN Bus.
    \item \textbf{ISO 15765(CAN-TP)}: el protocolo de transporte. El CAN Bus estándar solo envía 8 bytes por mensaje, pero la ECU necesita de un protocolo de transporte que le permita enviar mensajes de más de 8 bytes para poder enviar por ejemplo, el VIN del vehículo(17 bytes), este protocolo le permite fragmentar los mensajes en varios mensajes de 8 bytes y luego reensamblarlos en el destino.
\end{itemize}

\subsection{Diseño}

Para el diseño del projecto se

Ajuste a la planificación inicialmente prevista.

Modificación en los objetivos planteados.
